\documentclass[11pt,a4paper,fontset=fandol]{ctexart}

\usepackage[margin=2.25cm]{geometry}
\usepackage{amsmath,amssymb,mathtools,bm}
\usepackage{booktabs,longtable,tabularx,array,multirow}
\usepackage{graphicx}
\usepackage{xcolor}
\usepackage{enumitem}
\usepackage{microtype}
\usepackage{hyperref}
\usepackage[nameinlink,noabbrev]{cleveref}

\hypersetup{
  colorlinks=true,
  linkcolor=blue!55!black,
  citecolor=green!35!black,
  urlcolor=blue!60!black,
  pdftitle={从光滑近 Bessel 初值到两尺度轮廓},
  pdfauthor={IPM research note}
}
\graphicspath{{../../../result/verification/}}
\setlength{\emergencystretch}{3em}
\setlist{nosep,leftmargin=2em}
\renewcommand{\arraystretch}{1.16}

\newcommand{\R}{\mathbb{R}}
\newcommand{\dd}{\,\mathrm d}
\newcommand{\e}{\mathrm e}
\newcommand{\grad}{\nabla}
\newcommand{\perpgrad}{\nabla^\perp}
\newcommand{\abs}[1]{\left\lvert #1\right\rvert}
\newcommand{\norm}[1]{\left\lVert #1\right\rVert}
\newcommand{\Order}{\mathcal O}
\newcommand{\known}{\textcolor{green!45!black}{\textbf{严格/精确}}}
\newcommand{\asymstatus}{\textcolor{blue!65!black}{\textbf{渐近}}}
\newcommand{\numeric}{\textcolor{orange!70!black}{\textbf{本次数值}}}
\newcommand{\openq}{\textcolor{red!65!black}{\textbf{未知/猜想}}}

\newcommand{\resultfigure}[2][]{%
  \IfFileExists{../../../result/verification/#2}{%
    \includegraphics[#1]{../../../result/verification/#2}%
  }{%
    \fbox{\parbox[c][3.8cm][c]{0.85\linewidth}{\centering
      数值图尚未生成：\texttt{\detokenize{#2}}}}%
  }%
}

\InputIfFileExists{bessel_two_scale_metrics.tex}{}{%
  \newcommand{\NumFinalTime}{--}
  \newcommand{\NumFitMuInitial}{--}
  \newcommand{\NumFitMuFinal}{--}
  \newcommand{\NumFitHInitial}{--}
  \newcommand{\NumFitHFinal}{--}
  \newcommand{\NumProfileErrorInitial}{--}
  \newcommand{\NumProfileErrorFinal}{--}
  \newcommand{\NumKernelErrorInitial}{--}
  \newcommand{\NumKernelErrorFinal}{--}
  \newcommand{\NumWidthExponentInitial}{--}
  \newcommand{\NumWidthExponentFinal}{--}
  \newcommand{\NumZeroErrorInitial}{--}
  \newcommand{\NumZeroErrorFinal}{--}
  \newcommand{\NumGradientRatio}{--}
  \newcommand{\NumFieldChange}{--}
  \newcommand{\NumTailChange}{--}
  \newcommand{\NumRefinementDifference}{--}
  \newcommand{\NumKZeroPairFullRatio}{--}
  \newcommand{\NumKZeroPairTailRatio}{--}
  \newcommand{\NumKOneFitMuInitial}{--}
  \newcommand{\NumKOneFitMuFinal}{--}
  \newcommand{\NumKOneProfileErrorInitial}{--}
  \newcommand{\NumKOneProfileErrorFinal}{--}
  \newcommand{\NumKOneKernelErrorInitial}{--}
  \newcommand{\NumKOneKernelErrorFinal}{--}
  \newcommand{\NumKOneWidthExponentInitial}{--}
  \newcommand{\NumKOneWidthExponentFinal}{--}
  \newcommand{\NumKOneZeroErrorInitial}{--}
  \newcommand{\NumKOneZeroErrorFinal}{--}
  \newcommand{\NumKOneGradientRatio}{--}
  \newcommand{\NumKOneFieldChange}{--}
  \newcommand{\NumKOneTailChange}{--}
  \newcommand{\NumKOnePairFullRatio}{--}
  \newcommand{\NumKOnePairTailRatio}{--}
  \newcommand{\NumKOneWallLeakage}{--}
  \newcommand{\NumKOneExteriorError}{--}
  \newcommand{\NumKOneStopReason}{--}
  \newcommand{\NumStopReason}{--}
}

\title{\bfseries 从光滑近 Bessel 初值到两尺度轮廓\\
\large 精确结构、动态缩放、数值检验与已知/未知边界}
\author{IPM 研究报告}
\date{2026 年 8 月 28 日}

\begin{document}
\maketitle

\begin{abstract}
本文整理二维无粘不可压多孔介质方程（IPM）中右向聚焦 Bessel 解与
两尺度极限之间目前可以严格证明的联系，并把它与“光滑初值在真实时间
演化中自动收缩为两尺度奇异结构”这一尚未解决的问题分开。核心结论分为
Bessel 基准与全方向 $\rho_x\to0$ 的自屏蔽锐角候选两条主线。
第一，关系 $\rho=\lambda\psi$ 给出一族精确物理稳态；$K_0$ 成员在允许
壁面通量时具有非零壁面迹，$K_1\sin\theta$ 成员满足齐次不可穿壁条件，
二者均在右侧形成角宽 $\theta=\Order(r^{-1/2})$ 的抛物尾部。
第二，适当各向异性缩放后，$K_0$ 收敛到热核
$X^{-1/2}\exp(-\eta^2/(4X))$，而齐次壁面的 $K_1\sin\theta$ 收敛到
Dirichlet 热核的一阶法向形状
$\eta X^{-3/2}\exp(-\eta^2/(4X))$；它们确实是退化两尺度方程的稳态。
但对固定 Bessel 参数，这个极限对应向右远场扩展观察窗，而不是向物理尖点
收缩。第三，本文用当前不可穿壁上半平面求解器，对带小扰动的全局光滑近
$K_0$ 外极点初值和光滑正则化 $K_1$-core 初值进行短时演化，分别测量
密度、$\rho_x$ 零线、抛物宽度律和热核拟合误差。该实验是动力选择的
基准测试，不把初值自带的 Bessel 尾部误报成吸引子。最后给出真正收缩
尖扇所需的缩放关系、可检验量和下一步计算方案。第四，后续的 Poisson
多极矩审计修正了原三层解释：在
$\ell\ll r\ll\mu^{-1}$ 中，有界内核所需的 closure 修正必然大于
small-argument Bessel 分量，因此后者不能是主导中间层。Bessel 的正确角色是
$r\gtrsim\mu^{-1}$ 的外响应与漂移 Helmholtz Green 核；真正的多尺度候选还必须
包含 leading dynamic closure/wall layer 与全局 source-neutral return layer。
在额外硬条件 $\rho_x\in C_0$ 下，本文进一步构造了任意指定右锐扇内的
compact-angular leading tail，证明其高阶展开一般强制出现对数共振；又用
半平面 Green 逆矩公式构造了可自生正确双曲应变的光滑紧支撑 signed head。
散度恒等式排除了 finite-level-area 的静态 $\beta=0$ profile，壁面 Hilbert
闭合则在静态 compact edge 处强制平方根奇性。因此从光滑初值出发的最短
候选不是单一稳态，而是 local signed head、动态 fast edge 和固定物理尺度
Casimir cap 的非定常三层 gluing；其全局存在性仍是开放问题。
\end{abstract}

\begin{quote}
\textbf{2026-08-27 数学更正。}
本报告后文中的
$q_0\sim1/\log(1/\ell)$、$q_1\sim\ell$ 仅保留为历史性 amplitude matching，
不再是 source-free gluing 候选律。严格的单极/偶极矩排除、修正后的固定点方程与
一般未知 $0<\beta<1$ 的条件缩放律，以及后来发现的精确 affine 双曲内核与
$\theta_*=\pi\beta/(1+\beta)$ leading terminal fan，见
\texttt{analysis/BESSEL\_MULTISCALE\_BLOWUP\_THEORY\_ZH.md}。
先前代入的数值指数没有可靠数据或可解性依据，现已丢弃；主动理论只保留
一般未知的 $0<\beta<1$。
\end{quote}

\tableofcontents

\section{问题、约定与最短答案}

本文采用与仓库求解器一致的符号
\begin{equation}
  \rho_t+u\cdot\grad\rho=0,
  \qquad
  -\Delta\psi=\partial_x\rho,
  \qquad
  u=\perpgrad\psi=(-\psi_y,\psi_x).
  \label{eq:ipm}
\end{equation}
在上半平面 $\R\times\R_+$ 的不可穿壁问题中，壁面条件是
$u_2=\psi_x=0$；选定流函数规范后可写成 $\psi(x,0)=0$。
密度本身不需要在壁面为零。

针对“能否从光滑初值得到这个两尺度解”，截至本文日期，答案必须分成四层：
\begin{enumerate}
  \item \known：存在穿孔域精确/带 tip 分布源的奇异 Bessel 稳态核，并可严格
  推出两个 Gaussian 型两尺度极限；它们不是全域无源有限能稳态。
  \item \known：把 $K_0$ 极点放在壁面下方，可得到在闭上半平面内光滑的
  精确解；在一个\emph{扩展的}各向异性坐标系中，它从单个光滑状态严格
  收敛到 Gaussian 两尺度稳态。不过其解析速度穿过壁面。
  \item \openq：对无外力、不可穿壁的二维 IPM，尚无定理证明一般光滑
  初值在真实时间中自动收缩并选择这里的 Bessel/Gaussian 两尺度轮廓；
  本文数值只能提供短时筛查，不能填补这一逻辑缺口。
  \item \asymstatus：把远场缩放反向推进会依次看到 Gaussian 外区、
  Bessel 过渡层和中心核。这给出一个自然的 gluing 框架，但只有最后一层
  的收缩率由 IPM 动力学自身决定时，它才是物理时间中的 blow-up，而不只是
  对同一个静态场改变观察倍率。
\end{enumerate}

这一区分也解释了“它是不是 dynamic scaling 下的稳态”：
原始 Bessel 解在通常单尺度缩放下只有零缩放率时才是固定点；Gaussian
轮廓则在一个奇异的两尺度极限系统中是真固定点。后者对固定参数描述远场，
并不自动等价于有限时刻奇性。

\section{穿孔域精确/带 tip 源的 Bessel 稳态族}

\subsection{从 IPM 到修正 Helmholtz 方程}

令稳态密度与流函数线性相关：
\begin{equation}
  \rho=\lambda\psi,\qquad \lambda\ne0.
  \label{eq:rho-lambda-psi}
\end{equation}
于是 $u\cdot\grad\rho=\lambda\perpgrad\psi\cdot\grad\psi=0$，输运方程
自动成立，而 Poisson 方程化为
\begin{equation}
  \Delta\psi+\lambda\partial_x\psi=0.
  \label{eq:drift-helmholtz}
\end{equation}
作共轭变换
\begin{equation}
  \psi(x,y)=\e^{-\lambda x/2}\Phi(x,y),
  \qquad k=\frac{\abs\lambda}{2},
\end{equation}
便得到
\begin{equation}
  (\Delta-k^2)\Phi=0.
\end{equation}
因此在去掉极点的平面或扇形内，任意实阶 $\nu\ge0$ 的分离变量解
\begin{align}
  \psi_\nu(r,\theta)
  &=\e^{-\lambda r\cos\theta/2}K_\nu(kr)
    \bigl[a_\nu\cos(\nu\theta)+b_\nu\sin(\nu\theta)\bigr],
  \\
  \rho_\nu&=\lambda\psi_\nu
  \label{eq:bessel-family}
\end{align}
都是精确物理稳态。同一 $\lambda$ 下的收敛模态叠加仍然精确。

\subsection{右向聚焦的 K0 与不可穿壁的 K1}

取 $\lambda=-\mu<0$、$A>0$。全局右向聚焦的正密度成员为
\begin{equation}
  \psi_0=-A\e^{\mu x/2}K_0\!\left(\frac{\mu r}{2}\right),
  \qquad
  \rho_0=\mu A\e^{\mu x/2}K_0\!\left(\frac{\mu r}{2}\right).
  \label{eq:k0}
\end{equation}
它在穿孔平面精确，限制到上半平面后壁面流函数不为常数，因而是
给定通量/穿壁型成员，而非当前数值求解器的不可穿壁稳态。
若把它看成全平面分布解，则原点并非无源：
\begin{equation}
  (\Delta-\mu\partial_x)\psi_0=2\pi A\delta_0,
  \qquad
  \partial_x\rho_0-\mu^{-1}\Delta\rho_0=2\pi A\delta_0.
  \label{eq:k0-point-source}
\end{equation}
所以“全局”在这里意为穿孔平面的经典解及带点源的分布解，不能称为
无源全平面经典稳态。

满足 $\psi=0$ 于两条水平壁面射线的最低阶成员是
\begin{equation}
  \psi_1=-A\e^{\mu x/2}K_1\!\left(\frac{\mu r}{2}\right)\sin\theta,
  \qquad
  \rho_1=\mu A\e^{\mu x/2}K_1\!\left(\frac{\mu r}{2}\right)\sin\theta.
  \label{eq:k1}
\end{equation}
它符合不可穿壁边界，但在壁面点 $r=0$ 有偶极型奇性。
它还与 $K_0$ 成员满足精确恒等式
\begin{equation}
  \rho_1=-\frac{2}{\mu}\partial_y\rho_0.
  \label{eq:k1-derivative}
\end{equation}

利用 $K_\nu(z)\sim\sqrt{\pi/(2z)}\e^{-z}$~\cite{dlmf}，固定
$0<\theta<\pi$ 时
\begin{equation}
  \rho_\nu(r,\theta)
  \sim C_\nu(\theta)r^{-1/2}
  \exp\!\left[-\frac{\mu r}{2}(1-\cos\theta)\right].
  \label{eq:bessel-far}
\end{equation}
而在右射线附近
\begin{equation}
  1-\cos\theta=\frac{\theta^2}{2}+\Order(\theta^4),
  \qquad
  \rho_0\sim A\sqrt{\frac{\pi\mu}{r}}
  \exp\!\left(-\frac{\mu r\theta^2}{4}\right).
  \label{eq:parabolic-sector}
\end{equation}
故绝对横向宽度为 $y=\Order(\sqrt{r/\mu})$，相对角宽为
$\theta=\Order((\mu r)^{-1/2})$。它是全局解的缩角尾部，并非把扇形外
硬截为零；后者会在界面产生分布源。

\begin{figure}[htbp]
  \centering
  \resultfigure[width=0.93\linewidth]{exact_bessel_halfplane.png}
  \caption{右向聚焦 $K_0$ 穿壁成员与 $K_1\sin\theta$ 不可穿壁成员。
  图由 \texttt{ipm\_plot\_exact\_bessel\_family} 直接计算解析式。}
  \label{fig:exact-halfplane}
\end{figure}

\subsection{把极点移到壁面外：一个真正光滑但穿壁的例子}

对 $h>0$ 定义
\begin{equation}
  r_h=\sqrt{x^2+(y+h)^2},
  \qquad
  \psi_h=-A\e^{\mu x/2}K_0\!\left(\frac{\mu r_h}{2}\right),
  \qquad
  \rho_h=\mu A\e^{\mu x/2}K_0\!\left(\frac{\mu r_h}{2}\right).
  \label{eq:external-pole}
\end{equation}
极点位于 $(0,-h)$，所以 $(\rho_h,\psi_h)$ 在闭上半平面内解析、满足
内部方程、并在所有方向衰减到零。它给出了“单个光滑近 Bessel 状态”的
最干净构造。但
\begin{equation}
  u_2(x,0)=\partial_x\psi_h(x,0)\not\equiv0,
\end{equation}
所以它不属于不可穿壁半平面问题。该差别不是小边界误差：在适当衰减和
有限耗散条件下，不可穿壁平滑稳态由能量恒等式受到强刚性限制；有界域中
的 $C^1$ 稳态分类也只留下分层状态~\cite{elgindi2017}。

\section{两个严格的抛物两尺度极限}

\subsection{K0：热核极限}

为把大参数与远场缩放写在同一公式中，先令 Bessel 参数 $\mu=\varepsilon$，
并在 $X>0$ 上定义
\begin{equation}
  \mathcal R^{(0)}_\varepsilon(X,\eta)
  :=\varepsilon^{-1/2}
  \rho^{(0)}_\varepsilon\!\left(X,\frac{\eta}{\sqrt\varepsilon}\right)
  =A\sqrt\varepsilon\,\e^{\varepsilon X/2}
  K_0\!\left(
    \frac{\varepsilon}{2}\sqrt{X^2+\frac{\eta^2}{\varepsilon}}
  \right).
  \label{eq:k0-rescaled}
\end{equation}
在任意紧集 $X\in[X_0,X_1]$、$\abs\eta\le M$（$X_0>0$）上一致地
\begin{equation}
  \mathcal R^{(0)}_\varepsilon(X,\eta)
  \longrightarrow
  \mathcal R^{(0)}_*(X,\eta)
  :=A\sqrt{\frac{\pi}{X}}
  \exp\!\left(-\frac{\eta^2}{4X}\right).
  \label{eq:k0-limit}
\end{equation}
在 $X<0$ 的紧集上极限为零；跨过 $X=0$ 后，正确的分布表述是
\begin{equation}
  (\partial_X-\partial_{\eta\eta})\mathcal R^{(0)}_*
  =2\pi A\,\delta_{(0,0)}.
  \label{eq:heat-source}
\end{equation}
所以它是以下游空间 $X$ 为“热时间”的因果 Green 核，不是跨
$X=0$ 的光滑全局轮廓，更不是物理时间扩散。
直接展开还给出首项误差
\begin{equation}
  \mathcal R^{(0)}_\varepsilon
  =\mathcal R^{(0)}_*
  \left[1+\frac1\varepsilon\left(
  \frac{\eta^4}{16X^3}-\frac{\eta^2}{4X^2}-\frac1{4X}
  \right)+\Order(\varepsilon^{-2})\right].
  \label{eq:k0-correction}
\end{equation}
极限满足
\begin{equation}
  \partial_X\mathcal R^{(0)}_*
  =\partial_{\eta\eta}\mathcal R^{(0)}_*,
  \qquad
  X\mathcal R^{(0)}_{*,X}
  +\frac12\eta\mathcal R^{(0)}_{*,\eta}
  =-\frac12\mathcal R^{(0)}_*.
  \label{eq:k0-heat-homogeneity}
\end{equation}
横向质量和方差分别为
\begin{equation}
  \int_\R\mathcal R^{(0)}_*\dd\eta=2\pi A,
  \qquad
  \frac{\int_\R\eta^2\mathcal R^{(0)}_*\dd\eta}
       {\int_\R\mathcal R^{(0)}_*\dd\eta}=2X.
\end{equation}
事实上下游横向质量在有限 $\varepsilon$ 就精确守恒。由
\eqref{eq:k0-point-source} 可得原变量中
\begin{equation}
  \int_\R\rho_0(x,y)\dd y=
  \begin{cases}
    2\pi A, & x>0,\\
    2\pi A\e^{\mu x}, & x<0,
  \end{cases}
\end{equation}
而 $C_\rho C_y=1$，故对所有 $\varepsilon$ 和 $X>0$，
$\int_\R\mathcal R^{(0)}_\varepsilon\dd\eta=2\pi A$。这可作为解析与
MATLAB 两侧都很敏感的归一化检查。

用户特别关心的横向图像之外，$\rho_x$ 的符号结构也完全显式：
\begin{equation}
  \partial_X\mathcal R^{(0)}_*
  =\left(\frac{\eta^2}{4X^2}-\frac1{2X}\right)
  \mathcal R^{(0)}_*,
  \qquad
  \eta_{0}(X)=\sqrt{2X}.
  \label{eq:k0-rhox}
\end{equation}
因此零线以下为负、以上为正；有限 $\varepsilon$ 图中看到的弯曲换号线不是
数值噪声。由~\eqref{eq:k0-correction} 进一步得到 centered 分支的首修正
\begin{equation}
  \eta_0^2=2X-\frac{2}{\varepsilon}+\Order(\varepsilon^{-2}),
  \qquad X\ge X_0>0.
  \label{eq:k0-zero-correction}
\end{equation}

\subsection{K1 正弦模态：Dirichlet 热核的一阶形状}

对~\eqref{eq:k1} 取 $\mu=\varepsilon$，定义
\begin{equation}
  \mathcal R^{(1)}_\varepsilon(X,\eta)
  :=\rho^{(1)}_\varepsilon\!\left(X,\frac{\eta}{\sqrt\varepsilon}\right).
\end{equation}
由~\eqref{eq:k1-derivative}，在有限 $\varepsilon$ 就有
\begin{equation}
  \mathcal R^{(1)}_\varepsilon
  =-2\partial_\eta\mathcal R^{(0)}_\varepsilon,
  \label{eq:k1-finite-derivative}
\end{equation}
而不只是极限时近似相等。
则在 $X>0,\eta\ge0$ 的紧集上一致有
\begin{equation}
  \mathcal R^{(1)}_\varepsilon
  \longrightarrow
  \mathcal R^{(1)}_*(X,\eta)
  :=A\sqrt\pi\,\eta X^{-3/2}
  \exp\!\left(-\frac{\eta^2}{4X}\right)
  =-2A\sqrt\pi\,\partial_\eta
  \left[X^{-1/2}\e^{-\eta^2/(4X)}\right].
  \label{eq:k1-limit}
\end{equation}
这也是热方程解，并在 $\eta=0$ 为零。其抛物齐次次数为 $-1$，且
\begin{equation}
  \partial_X\mathcal R^{(1)}_*
  =\left(\frac{\eta^2}{4X^2}-\frac{3}{2X}\right)
   \mathcal R^{(1)}_*,
  \qquad
  \eta_0(X)=\sqrt{6X}.
  \label{eq:k1-rhox}
\end{equation}
所以不可穿壁外解与穿壁 $K_0$ 解的两尺度目标不同：不能用同一个 Gaussian
模板评价两者。centered 有限参数分支的非平凡零线还有首修正
\begin{equation}
  \eta_0^2=6X-\frac{6}{\varepsilon}+\Order(\varepsilon^{-2}).
  \label{eq:k1-zero-correction}
\end{equation}
若在全平面作奇延拓，其分布源是
$-4\pi A\partial_\eta\delta_{(0,0)}$；在半平面则应理解为 tip 处的
边界偶极源。相应地，有限 $\varepsilon$ 的上半线第一矩精确为
\begin{equation}
  \int_0^\infty \eta\,\mathcal R^{(1)}_\varepsilon(X,\eta)\dd\eta
  =2\pi A,\qquad X>0.
  \label{eq:k1-first-moment}
\end{equation}

\begin{figure}[htbp]
  \centering
  \resultfigure[width=0.9\linewidth]{exact_bessel_halfplane_rho_x.png}
  \caption{精确 Bessel 半平面成员的 $\rho_x$。$K_0$ 与 $K_1$ 两类的
  抛物极限零线分别是 $\eta=\sqrt{2X}$ 和 $\eta=\sqrt{6X}$。}
  \label{fig:exact-rhox}
\end{figure}

\section{各向异性 dynamic scaling：哪里是真稳态，哪里只是坐标效应}

\subsection{完整两尺度变换}

令
\begin{equation}
  X=C_x(t)(x-a(t)),\qquad Y=C_y(t)y,\qquad
  R=C_\rho(t)\rho,\qquad
  P=\frac{C_\rho C_y^2}{C_x}\psi,
  \label{eq:anisotropic-map}
\end{equation}
并选择保留薄扇速度为 $\Order(1)$ 的时间
\begin{equation}
  \frac{\dd t}{\dd s}=\frac{C_\rho C_y}{C_x^2}.
\end{equation}
记
\begin{equation}
  \delta=\left(\frac{C_x}{C_y}\right)^2,\qquad
  c_x=\partial_s\log C_x,\quad
  c_y=\partial_s\log C_y,\quad
  c_\rho=\partial_s\log C_\rho,
\end{equation}
以及平移率
$c_r=-(C_\rho C_y/C_x)\dot a$。精确变换后的系统为
\begin{align}
  R_s+\bigl(-P_Y+c_xX+c_r\bigr)R_X
      +\bigl(P_X+c_yY\bigr)R_Y&=c_\rho R,
  \label{eq:two-scale-transport}\\
  -\bigl(\delta P_{XX}+P_{YY}\bigr)&=R_X,
  \label{eq:two-scale-poisson}\\
  \delta_s&=2(c_x-c_y)\delta.
  \label{eq:delta-evolution}
\end{align}
当 $C_x=C_y$ 时，这个时钟和流函数归一化恰好退回通常单尺度约定。
这就是将来两尺度求解器应实现的方程；当前仓库的 dynamic 模式只有
$C_x=C_y$，不能锁定两个独立宽度。

薄扇极限是 $\delta\to0$：
\begin{equation}
  -P_{YY}=R_X.
  \label{eq:degenerate-limit}
\end{equation}
注意，\eqref{eq:degenerate-limit} 本身不是热方程。对 Bessel 轨道，
$C_y^2=\mu C_x$ 且 $\rho=-\mu\psi$，因而有额外的精确关系
\begin{equation}
  P=-R.
  \label{eq:p-minus-r}
\end{equation}
只有把它代入椭圆方程，才得到
$R_X=R_{YY}+\delta R_{XX}$，并在极限得到热方程。速度
$(-P_Y,P_X)=(R_Y,-R_X)$ 仍为 $\Order(1)$，只是因为与 $\grad R$
正交而使 Bessel 非线性输运精确消失。因此 Gaussian 是两尺度边界系统中
一个特殊的固定点，但一般 IPM 极限并不会自动选择关系~\eqref{eq:p-minus-r}。
在这条特殊支上，极限速度也完全显式：
\begin{equation}
  V_*=(-P_{*,Y},P_{*,X})=(R_{*,Y},-R_{*,X}),
  \qquad V_*\cdot\grad R_*=0.
\end{equation}
对 $K_0$，
$V_{*,2}=(1/(2X)-Y^2/(4X^2))R_*^{(0)}$，故壁面仍有穿流；
对 $K_1$，$R_*^{(1)}(X,0)=R_{*,X}^{(1)}(X,0)=0$，故
$V_{*,2}|_{Y=0}=0$。Poisson、边界条件和输运几何由此同时闭合。

\subsection{一个严格的含时收敛轨道}

对固定物理参数 $\mu$，令
\begin{equation}
  C_x(s)=\mu\e^{-\gamma s},
  \qquad
  C_y(s)=\mu\e^{-\gamma s/2},
  \qquad
  C_\rho(s)=C_y^{-1},
  \qquad \gamma>0.
  \label{eq:expanding-gauge}
\end{equation}
于是
\begin{equation}
  \varepsilon=\frac{\mu}{C_x}=\e^{\gamma s}\to\infty,
  \qquad
  (c_x,c_y,c_\rho)=(-\gamma,-\gamma/2,\gamma/2).
\end{equation}
将一个不随物理时间变化的 $K_0$ 解写入这些坐标，
$R(s)\to\mathcal R^{(0)}_*$，由~\eqref{eq:k0-correction} 收敛率为
$\Order(\e^{-\gamma s})$。极限的齐次恒等式保证
\begin{equation}
  c_xX R_X+c_yYR_Y=c_\rho R.
\end{equation}
所以它确实是含时 dynamic scaling 下的新稳态。

但是~\eqref{eq:expanding-gauge} 对应的物理长度
$L_x=C_x^{-1}\sim\e^{\gamma s}$、
$L_y=C_y^{-1}\sim\e^{\gamma s/2}$ 都在增长。它在向右看越来越远，
而不是向一个点越缩越小。若采用通常的尖点放大 $C_x\to\infty$，固定
$\mu$ 反而给出 $\varepsilon=\mu/C_x\to0$，看到的是 $K_0$ 的对数核心，
不是 Gaussian 尾部。
各向异性缩放后的时钟并不唯一；改变时钟会改变 $c_x$ 的绝对数值和
$C_x$ 对物理时间的日程，但不会改变
$c_y/c_x=1/2$、$c_\rho/c_x=-1/2$ 或“两个物理长度都增长”这一几何事实。
在这条显式轨道上物理解本身完全不动，所谓收敛全部来自观察窗。
上述 $\Order(\e^{-\gamma s})$ 是极点居中的速率。若使用固定深度 $h>0$
的域外极点而不在 $Y$ 中重新定心，其缩放位移
$H=C_yh=\mu h\varepsilon^{-1/2}$，故首项误差一般降为
$\Order(\varepsilon^{-1/2})=\Order(\e^{-\gamma s/2})$；平移调制掉 $H$
后才恢复居中支的 $\Order(\varepsilon^{-1})$ 修正。

对 $K_1$ 极限应改用 $C_\rho=C_x^{-1}$，从而
$c_\rho=-c_x$，与次数 $-1$ 的抛物齐次律一致。

\subsection{反向缩放、中心核与三层 gluing}

上面的 blow-down 结论只描述 $\varepsilon\to\infty$。反向令
$C_x\to\infty$，则
\begin{equation}
  \varepsilon=\frac{\mu}{C_x}\to0,
  \qquad L_x=C_x^{-1}\to0,
\end{equation}
观察窗确实向中心收缩。因此从重整化尺度流的意义看，
\begin{equation}
  \text{Gaussian 远场}
  \quad\xrightarrow{\ \varepsilon\downarrow0\ }\quad
  \text{Bessel 中心核}
  \label{eq:gaussian-to-bessel-rg}
\end{equation}
是一个 blow-up 方向。不过，对固定 $\mu$ 的精确 Bessel 成员，这仍只是
沿同一个静态空间场由外向内扫描，并不是物理时间把一个全局 Gaussian
初值演化成新生奇点。

中心渐近式把缺少的内层结构显示出来。由
$K_0(z)=-\log(z/2)-\gamma_E+\Order(z^2\abs{\log z})$，
\begin{equation}
  \rho_0(r,\theta)
  =\mu A\left[-\log\frac{\mu r}{4}-\gamma_E\right]
  +\Order\!\left(\mu^2r\abs{\log(\mu r)}\right),
  \qquad
  \abs{\grad\rho_0}\sim\frac{\mu A}{r}.
  \label{eq:k0-core}
\end{equation}
同样，由 $K_1(z)=z^{-1}+\Order(z\abs{\log z})$，
\begin{equation}
  \rho_1(r,\theta)
  =\frac{2A\sin\theta}{r}+\Order\!\left(\mu\right)
  =2A\frac{y}{r^2}+\Order\!\left(\mu\right),
  \qquad
  \abs{\grad\rho_1}=\Order(r^{-2}).
  \label{eq:k1-core}
\end{equation}
这说明完整候选结构不应只分成 Gaussian 与 Bessel 两层，而应写成
\begin{equation}
  \underbrace{r\sim\ell(t)}_{\text{有界动态内核}}
  \quad\longleftrightarrow\quad
  \underbrace{\ell(t)\ll r\lesssim\mu^{-1}}_{\text{Bessel 中间区}}
  \quad\longleftrightarrow\quad
  \underbrace{x\gg\mu^{-1},\;
    y=\Order(\sqrt{x/\mu})}_{\text{Gaussian 双尺度外区}},
  \qquad \ell(t)\downarrow0.
  \label{eq:three-layer-gluing}
\end{equation}
外区与中间区的 gluing 是“自动”的：Gaussian 正是同一 Bessel 函数的
远场展开。真正非平凡的是内区与中间区的动态 gluing，因为光滑 IPM 输运
在经典解存在期间满足
\begin{equation}
  \norm{\rho(t)}_{L^\infty}=\norm{\rho(0)}_{L^\infty};
  \label{eq:linfty-transport}
\end{equation}
所以最终只能让梯度发散，不能直接形成无界的 $K_0$ 或 $K_1$ 密度核。

这一约束立即给出两条形式匹配律。对 $K_0$ 型中间区写
\begin{equation}
  \rho_{\rm mid}(t,r)
  \sim b(t)-q_0(t)\log r.
\end{equation}
为了让 $r=\ell(t)$ 处的密度振幅保持 $\Order(1)$，临界幅度必须满足
\begin{equation}
  q_0(t)\log\frac1{\ell(t)}=\Order(1),
  \qquad
  q_0(t)\sim\frac{1}{\log(1/\ell(t))}.
  \label{eq:k0-gluing-law}
\end{equation}
于是仍可有
\begin{equation}
  \abs{\grad\rho}_{r\sim\ell(t)}
  \sim\frac{1}{\ell(t)\log(1/\ell(t))}\to\infty.
  \label{eq:k0-glued-gradient}
\end{equation}
这表明先前考虑的对数修正可能不是次要装饰，而是由 $L^\infty$ 守恒和
内外匹配共同强制出来的。

对 $K_1$ 型偶极中间区写
\begin{equation}
  \rho_{\rm mid}(t,x,y)\sim q_1(t)\frac{y}{r^2}.
\end{equation}
在 $r=\ell(t)$ 处保持有界要求
\begin{equation}
  q_1(t)=\Order(\ell(t)),
  \qquad
  \abs{\grad\rho}_{r\sim\ell(t)}
  =\Order\!\left(\frac1{\ell(t)}\right)\to\infty.
  \label{eq:k1-gluing-law}
\end{equation}
该支天然具有正负偶极结构，因此还可能与零质量、第一法向矩抵消及锐角扇
的全局闭合更相容。

式~\eqref{eq:k0-gluing-law}--\eqref{eq:k1-gluing-law} 只是由匹配和最大值
约束得到的必要尺度律，不是存在性证明。严格构造还需把三层复合近似代回
IPM，对平移、幅度和缩放零模施加可解性条件，从而推出 $\ell(t)$、$q_j(t)$
及可能的 $\mu(t)$ 的调制方程。只有当该方程由真实时间动力学给出
$\ell(t)\downarrow0$，式~\eqref{eq:gaussian-to-bessel-rg} 才从静态
重整化轨道升级为真正的 blow-up 机制。

\subsection{全局质量约束}

若一个各向异性固定点足够快衰减并且边界通量消失，则
\begin{equation}
  \grad\cdot\left(R[U+(c_xX,c_yY)]\right)
  =(c_x+c_y+c_\rho)R
\end{equation}
给出必要条件
\begin{equation}
  (c_x+c_y+c_\rho)\int R\dd X\dd Y=0.
  \label{eq:mass-constraint}
\end{equation}
$K_0$ 热核沿 $X$ 方向具有常数横向质量，故总质量发散，
\eqref{eq:mass-constraint} 不适用。任何声称它是非负有限质量全局吸引子的
论证，都必须另加外区匹配、零总质量或改变幅度规范。

在不可穿壁半平面的薄扇极限，还有更贴切的一阶矩条件。若
$P(X,0)=0$ 且 $P,P_Y\to0$ 当 $Y\to\infty$，由
$-P_{YY}=R_X$ 两次积分得到
\begin{equation}
  \partial_X J(X)=0,\qquad
  J(X):=\int_0^\infty YR(X,Y)\dd Y.
  \label{eq:first-moment-constraint}
\end{equation}
一个在 $X$ 两端都局域、非负、无源且无外区调和势的轮廓只能有
$J=0$，从而只能平凡。K1 极限以~\eqref{eq:k1-first-moment} 的非零常数
$J=2\pi A$ 和 tip 偶极源规避这一障碍。这进一步说明：真正有限质量的
无源两尺度固定点必须包含外区匹配或不同符号结构。

\section{真正“收缩锐角扇”的两尺度候选}

Bessel 热核尾部是 $r\to\infty$ 的缩角结构。若目标是在壁面点附近
$x\downarrow0$ 时形成越来越窄的锐角扇，应直接写
\begin{equation}
  \eta=\frac{y}{x^p},\qquad p=1+\sigma>1,
  \qquad
  Q(x,y)=x^\beta F(\eta),
  \qquad
  \psi(x,y)=x^\gamma G(\eta).
  \label{eq:shrinking-fan}
\end{equation}
这对应绝对横向尺度 $L_y=L_x^p\ll L_x$，与 Bessel 远场的
$L_y\sim L_x^{1/2}$ 方向相反。

由于 $p>1$，Poisson 方程的法向二阶导数占优。配平
$-\psi_{yy}=Q_x$ 得到
\begin{equation}
  \gamma=\beta+2p-1=\beta+1+2\sigma,
  \qquad
  -G''=\beta F-p\eta F'.
  \label{eq:shrinking-poisson}
\end{equation}
非线性输运的首项出现一个有用的消去：
\begin{equation}
  u\cdot\grad Q
  =x^{2\beta+\sigma-1}
  \left[-\beta G'F+\gamma GF'\right]+\mathrm{l.o.t.}
  \label{eq:shrinking-transport}
\end{equation}
相对于缩放项 $x^\beta$，临界配平是
\begin{equation}
  \beta+\sigma=1.
  \label{eq:critical-fan-relation}
\end{equation}
若壁面 $Q_x\sim x^{-\alpha}$，则 $\beta=1-\alpha$，临界情形给出
$\sigma=\alpha$。这里不预设任何数值 $\alpha$；
\eqref{eq:critical-fan-relation} 是将来可从数值同时测量的关系，而非已知指数。

一般定常两尺度轮廓还必须满足角向方程
\begin{equation}
  -\beta G'F+\gamma GF'+(\beta c_x-c_\rho)F=0,
  \qquad c_y=p c_x,
  \label{eq:shrinking-angular-dynamic}
\end{equation}
而不是自动满足 $c_\rho=\beta c_x$。只有当首阶非线性输运单独消失时，
$\beta c_x=c_\rho$，此时
\begin{equation}
  F=C\abs{G}^{\beta/\gamma}.
  \label{eq:shrinking-streamline}
\end{equation}
齐次不可穿壁条件 $G(0)=0$ 又会迫使连续的 $F(0)=0$。因此两尺度本身
并没有自动解决“壁面上非零奇异迹”的兼容性问题。不过一般角向方程比这
个特殊情形更有空间：形式积分给出
\begin{equation}
  F=C\abs{G}^{\beta/\gamma}
  \exp\!\left[-\frac{\beta c_x-c_\rho}{\gamma}
  \int^\eta\frac{\dd z}{G(z)}\right].
  \label{eq:shrinking-general-F}
\end{equation}
若 $G(\eta)=a\eta+\cdots$ 且希望 $F(0)=F_0\ne0$，则必要的壁面配平是
\begin{equation}
  c_\rho=\beta(c_x-a),
  \qquad
  G''(0)=-\beta F_0,
  \qquad
  F'(0)=-\frac{\beta^2F_0^2}{\gamma a}\quad(a\ne0).
  \label{eq:shrinking-wall-compatibility}
\end{equation}
因此缩放源与壁面切向速度可以在形式上支持非零壁面迹；这给出了一个新的
shooting 边值问题，但尚不是存在性结论。无论采用哪一支，必要的几何率
关系只有 $c_y/c_x=p$。这比单一 $C_l$ 更直接：数值中应分别记录纵向
位置/长度和法向宽度，而不是事先锁死某个指数。

在薄扇时钟下，真正近尖点极限应直接求解
\begin{align}
  R_s+(-P_Y+c_xX)R_X+(P_X+(1+\sigma)c_xY)R_Y&=c_\rho R,\\
  -P_{YY}&=R_X,
  \label{eq:shrinking-limit-system}
\end{align}
并取 $c_x>0$。若还要求有限二维质量且缩放边界通量为零，
\eqref{eq:mass-constraint} 强制
\begin{equation}
  c_\rho=-(2+\sigma)c_x.
\end{equation}
它与“首阶输运单独消失”这条特殊支上的
$c_\rho=\beta c_x$（$\beta>0$）不相容，正说明这条特殊局部轮廓只能是
内区，必须与携带质量的外区匹配。一般支仍须把
\eqref{eq:shrinking-angular-dynamic}、质量通量和外区同时求解。
这与 Bessel blow-down 的 $c_y/c_x=1/2$、$c_\rho/c_x=-1/2$
完全不同。把 Bessel 的 rate 直接搬到收缩尖扇会违反质量与几何方向。

\section{角度依赖的指数是否有帮助}

考虑更一般的 WKB 形式
\begin{equation}
  \rho(r,\theta)=f(\theta)r^{\beta(\theta)}
  \exp\!\left[g(\theta)r^{\alpha(\theta)}\right].
  \label{eq:angle-dependent-alpha}
\end{equation}
角导数指数部分包含
\begin{equation}
  \partial_\theta\bigl(g r^\alpha\bigr)
  =r^\alpha\left[g'+g\alpha'\log r\right].
\end{equation}
在 $g\ne0$ 的开区间，只要 $\alpha'\ne0$，角向二阶项会产生
$r^{2\alpha-2}(\alpha'\log r)^2$，比没有对数的径向 WKB 项更强。
对常系数的修正 Helmholtz 方程，这迫使首阶通常满足 $\alpha'=0$。
因此：
\begin{itemize}
  \item 在一个固定开扇形内部，先让 $\alpha$ 为常数、把方向依赖放进
  $g(\theta)$ 更自然；Bessel 家族正是 $\alpha=1$、
  $g=-\mu(1-\cos\theta)/2$。
  \item $\alpha(\theta)$ 仍可能在 $g=0$ 的转向射线、分片 WKB 匹配区或
  shrinking-fan 边界层中有用，但必须显式保留 $\log r$ 项，不能把它当成
  普通缓变参数。
\end{itemize}

\section{从光滑初值到两尺度：文献中已知什么}

本节只讨论与当前问题直接相关的结果。截至 2026 年 8 月 27 日：

\begin{longtable}{>{\raggedright\arraybackslash}p{0.20\linewidth}
                  >{\raggedright\arraybackslash}p{0.34\linewidth}
                  >{\raggedright\arraybackslash}p{0.37\linewidth}}
\toprule
结果 & 已证明内容 & 对本文问题的限制 \\
\midrule
\endfirsthead
\toprule
结果 & 已证明内容 & 对本文问题的限制 \\
\midrule
\endhead
Castro--C\'ordoba--Gancedo--Orive~\cite{castro2007}
& 光滑解的局部存在唯一性、延拓判据；特殊无限能类中有显式爆破机制。
& 不给出一般二维光滑有限能无外力数据到 Bessel 两尺度的选择。\\

Elgindi~\cite{elgindi2017}
& 在适当条件下分类稳态，并证明稳定分层稳态附近的小扰动全局存在和渐近稳定。
& 支持“平滑、衰减、不可穿壁的非分层 Bessel 型物理稳态受到阻碍”；本文
带 tip 源/非标准远场的核不属于该经典稳态类。\\

Bianchini 等~\cite{bianchinijoparkwang2025,bianchinicordoba2025}
& 锐化稳定分层态附近的渐近稳定性；另一方面证明临界 $H^2$ 中强不适定。
& 表明光滑相空间同时存在稳定吸引域和剧烈不稳定阈值，不可能由所有光滑
初值普遍选择同一个集中轮廓。\\

Kiselev--Yao~\cite{kiselevyao2023}
& 构造数据：若解始终光滑，则其导数必须在无限时间中无界增长；证明某些分层
稳态的非线性不稳定。
& 是小尺度形成的严格证据，但不证明有限时爆破，也不识别本文两尺度轮廓。\\

Kiselev--Sarsam~\cite{kiselevsarsam2025}
& 从半平面边界层推导一维非局部 IPM 模型，并对光滑有界周期数据证明有限时爆破。
& 是一维模型定理，不是二维 IPM 定理。\\

C\'ordoba--Mart\'inez-Zoroa~\cite{cordobamartinez2024}
& 对带光滑紧支撑外源的二维 IPM，从光滑紧支撑初值构造有限时 $C^1$ 奇性。
& 方程含外源 $F$；不能移除 $F$ 后直接用于本文无外力问题。\\

Collot--Prange--Tan~\cite{collotprangetan2026}
& 对一个特殊无限能降维类中的显式自相似爆破解证明光滑扰动稳定性。
& 稳定性位于特殊 ansatz/无限能类内，目标轮廓也不是 Bessel 两尺度。\\

Dembski~\cite{dembski2025}
& 在严格小于半平面的锐角域，对原点处仅 Lipschitz、角变量高阶光滑且可在
远处截断的数据证明无外力 IPM 有限时奇性，并给出有限余维稳定构造。
& 初值在角点不是 $C^1$，域是锐角而非光滑全平面/标准半平面；仍不回答
一般二维 $C^\infty$ 初值的选择问题。\\

Wang 等~\cite{wang2025,wangleger2025}
& 在不可穿壁上半平面数值发现多个 dynamic-scaling 稳态候选，并以高精度
残差和数值线性化区分稳定/不稳定模态。
& 目前是浮点/PINN 数值证据，且使用各向同性单尺度，不是本文的薄扇
two-scale 极限。\\
\bottomrule
\end{longtable}

综合这些结果，下面这句话目前仍然是开放的：
\begin{quote}
存在一个无外力、不可穿壁的二维 IPM 光滑有限能初值，其真实时间解在有限
或无限时间中，经两个独立收缩尺度归一化后，收敛到本文的 Bessel/Gaussian
轮廓。
\end{quote}
特别地，“已有平滑爆破解”“已有稳定自相似解”和“已有锐角 Lipschitz
爆破解”都不能替代这条尚缺的桥梁。
全平面还存在沿有限宽斜条带支撑的光滑稳态~\cite{constantinlavicol2019}，
但它们沿条带轴向不局域，也不提供右向 Bessel 尾部的动力选择。

\section{MATLAB：光滑近 Bessel 初值的动力测试}

\subsection{两类初值与边界契约}

第一类是外极点 K0 数据
\begin{equation}
  \rho^{K0}_{0,h,\delta}(x,y)
  =C_h\e^{\mu x/2}K_0\!\left(
  \frac\mu2\sqrt{x^2+(y+h)^2}\right)
  \left[1+\delta B(x,y)\right],
  \label{eq:numerical-initial}
\end{equation}
其中 $C_h$ 令网格最大值为 $1$，$B$ 是在壁面消失的光滑局域 bump。
采用 $\mu=2,h=0.6,\delta=0.03$。极点在计算域外，所以密度初值
全局光滑且向远端衰减。加入 $\delta\ne0$ 是必要的：只画无扰动 Bessel
尾部只能验证渐近公式，不能测试吸引性。

第二类是更贴合不可穿壁外解的光滑 K1-core 数据。令
\begin{equation}
  s_{\epsilon_c}=\sqrt{x^2+y^2+\epsilon_c^2},\qquad
  \rho^{K1}_{0,\epsilon_c,\delta}
  =C_1\e^{\mu x/2}K_1\!\left(\frac{\mu s_{\epsilon_c}}2\right)
  \frac{y}{s_{\epsilon_c}}\,[1+\delta B(x,y)],
  \label{eq:k1-core-initial}
\end{equation}
其中 $\epsilon_c=0.45$。纯 K1-core 数据是 $C^\infty$ 的并在 $y=0$
精确为零；独立验证得到壁面密度与法向速度均为机器零，$r\ge4$ 外区对精确
K1 的相对误差为 \NumKOneExteriorError。由于当前通用求解器的物理模式
安全审计仍要求一个正的壁面 $\rho_x$ 峰，实际 K1 演化另加了振幅
$5\times10^{-3}$ 的 signed-Gaussian wall monitor；其末态壁面泄漏占峰值
\NumKOneWallLeakage。该 monitor 是已报告的数值兼容项，不属于 K1 外解。

当前求解器在 $[-24,24]\times[0,12]$ 上强制 $\psi=0$ 于 $y=0$，
人工远边界使用 half-plane Green 数据，密度使用开放输运边界。对 K0，
它会从~\eqref{eq:numerical-initial} 重新求解不可穿壁速度，而不是沿解析
穿壁 $(\rho_h,\psi_h)$ 原地不动；对 K1-core，它同样重算流函数，光滑核心
使精确稳态关系只在外区近似成立。

两类主筛查均使用 $257\times129$ 非均匀网格、WENO5、SSPRK3、
$\Delta t_{\max}=2.5\times10^{-4}$，共 200 步演化到
$t=\NumFinalTime$。本轮尚未完成同终止时刻的跨网格收敛研究，因此所有
动力结论标为筛查结果。运行入口仍为唯一的 \texttt{main\_ipm}；可复现
选项在 \path{experiments/bessel_two_scale_case.m}，K0/K1 分析分别在
\path{analysis/ipm_analyze_bessel_two_scale.m} 与
\path{analysis/ipm_analyze_bessel_k1_two_scale.m}。

\subsection{检验量}

对尾区每个 $x>0$，外极点 $K_0$ 的抛物预测为
\begin{equation}
  \rho(x,y)\approx Cx^{-1/2}
  \exp\!\left[-\frac{\mu(y+h)^2}{4x}\right].
  \label{eq:numerical-kernel}
\end{equation}
K1 对应模板为
\begin{equation}
  \rho(x,y)\approx C yx^{-3/2}
  \exp\!\left[-\frac{\mu y^2}{4x}\right],
  \qquad
  y_{\rm peak}=\sqrt{\frac{2x}{\mu}},
  \qquad
  y_0=\sqrt{\frac{6x}{\mu}}.
\end{equation}
两个分析程序不只拟合热图，而是独立计算：
\begin{enumerate}
  \item 由 $-x\log[\rho(x,y)/\rho(x,0)]$ 对 $y^2,y$ 的最小二乘拟合
  $\mu_{\rm fit},h_{\rm fit}$；
  \item 全热核相对 $L^2$ 误差与归一化横截面误差；
  \item $1/e$ 横向宽度对 $x$ 的对数斜率，预测值为 $1/2$；
  \item 直接差分所得 $\rho_x$ 的换号线，与
  \begin{equation}
    y_0(x)=\sqrt{\frac{2x}{\mu}}-h
  \end{equation}
  比较；
  \item 全场/尾区初末变化以及 $\norm{\grad\rho}_\infty$ 增长比，用来判断
  “只是保留初值尾形”还是出现了实际动力学选择。
\end{enumerate}

\subsection{结果}

\begin{table}[htbp]
\centering
\caption{带 $3\%$ 光滑扰动的 K0 外极点主筛查。误差均为相对 $L^2$ 量。}
\label{tab:numerical-results}
\begin{tabular}{lcc}
\toprule
量 & $t=0$ & $t=\NumFinalTime$ \\
\midrule
$\mu_{\rm fit}$ & \NumFitMuInitial & \NumFitMuFinal \\
$h_{\rm fit}$ & \NumFitHInitial & \NumFitHFinal \\
归一化横截面误差 & \NumProfileErrorInitial & \NumProfileErrorFinal \\
全热核误差 & \NumKernelErrorInitial & \NumKernelErrorFinal \\
宽度指数 & \NumWidthExponentInitial & \NumWidthExponentFinal \\
$\rho_x$ 零线误差 & \NumZeroErrorInitial & \NumZeroErrorFinal \\
\bottomrule
\end{tabular}
\end{table}

全场初末相对变化为 \NumFieldChange，尾区相对变化为 \NumTailChange，
$\norm{\grad\rho}_\infty$ 增长比为 \NumGradientRatio。匹配的
扰动/无扰动两条轨道之差，在全场和尾区的末初比值分别为
\NumKZeroPairFullRatio、\NumKZeroPairTailRatio；求解器停止原因为
\texttt{\NumStopReason}。

\begin{figure}[htbp]
  \centering
  \resultfigure[width=\linewidth]{bessel_two_scale_perturbed_diagnostics.png}
  \caption{平滑近 $K_0$ 初值在不可穿壁 IPM 下的短时演化。上排依次为
  初始密度、末态密度、末态 $\rho_x$ 及解析零线；下排给出拟合参数、
  三类缺陷和抛物横截面 collapse。}
  \label{fig:numerical-diagnostics}
\end{figure}

\begin{table}[htbp]
\centering
\caption{带 $3\%$ 光滑扰动和 $0.5\%$ signed-Gaussian wall monitor 的
K1-core 主筛查。}
\label{tab:k1-numerical-results}
\begin{tabular}{lcc}
\toprule
量 & $t=0$ & $t=\NumFinalTime$ \\
\midrule
$\mu_{\rm fit}$ & \NumKOneFitMuInitial & \NumKOneFitMuFinal \\
归一化横截面误差 & \NumKOneProfileErrorInitial & \NumKOneProfileErrorFinal \\
全热核导数误差 & \NumKOneKernelErrorInitial & \NumKOneKernelErrorFinal \\
峰位置宽度指数 & \NumKOneWidthExponentInitial & \NumKOneWidthExponentFinal \\
$\rho_x$ 零线误差 & \NumKOneZeroErrorInitial & \NumKOneZeroErrorFinal \\
\bottomrule
\end{tabular}
\end{table}

K1 的全场/尾区初末相对变化为
\NumKOneFieldChange、\NumKOneTailChange，梯度增长比为
\NumKOneGradientRatio；匹配轨道差的全场/尾区末初比为
\NumKOnePairFullRatio、\NumKOnePairTailRatio，停止原因为
\texttt{\NumKOneStopReason}。

\begin{figure}[htbp]
  \centering
  \resultfigure[width=\linewidth]{bessel_k1_two_scale_perturbed_diagnostics.png}
  \caption{光滑 K1-core 初值的短时演化。末态 $\rho_x$ 黑色虚线使用
  Dirichlet 热核导数预测 $y^2=6x/\mu_{\rm fit}$。wall monitor 的泄漏
  单独记录，没有被当作 K1 主轮廓。}
  \label{fig:k1-numerical-diagnostics}
\end{figure}

两类数据都清楚显示短时保形，但没有显示吸引：K0 的匹配扰动差只缩小约
$0.15\%$--$0.18\%$，K1 的比值在 $1$ 附近且尾区略增。尾部拟合误差也只在
第四位变化。因此本轮结论是“预置 Bessel 尾部在短时间内保持”，不是
“任意光滑数据收敛到 Bessel 两尺度”。即使未来观察到误差下降，也仍须
通过更长时间、跨网格和无预置长尾测试，才能讨论吸引或有限时爆破。

\section{全方向 \texorpdfstring{$\rho_x$}{rho-x} 衰减后的结论更新}

现在把远场条件明确为二维一致衰减
\begin{equation}
  \rho_x(t,\cdot)\in C_0(\overline{\mathbb H}),
  \qquad
  \lim_{|(x,y)|\to\infty,\,y\geq0}\rho_x(t,x,y)=0,
  \label{eq:rhox-c0}
\end{equation}
并排除 $\rho_x\equiv0$、只有 $\rho_y$ 爆破的退化情形。这个条件立即淘汰
已有的水平 front、旋转无限条带和 affine 锐扇：它们分别在 $y$ 方向、条带
切向或所有空间点保持非零水平梯度。singular $K_0/K_1$ 的导数确实衰减，
但它们在每个时刻已经带 tip 奇点并且在 $D=0$ 类中静止，不能作为光滑初值
爆破。

Dembski~\cite{dembski2025} 的严格锐角域结果则可把密度在空间无穷远截成
紧支撑，因而满足~\eqref{eq:rhox-c0}，并确实形成有限时间梯度奇点。它的
初值在角点一般只到 Lipschitz，且定义在严格小于半平面的锐角域，所以没有
解决“全局光滑半平面初值”这一更强目标；但它已经证明远端衰减本身不阻止
IPM 爆破。

\subsection{\texorpdfstring{$\beta=0$}{beta=0} bounded-contrast 与非定常自屏蔽}

放弃 $R\in L^p$、但保留~\eqref{eq:rhox-c0} 后，最短的新候选取
\begin{equation}
  L(t)=c(T-t),\qquad
  \rho=\rho_c+R\!\left(\frac{x-a(t)}{L(t)},\frac{y}{L(t)}\right),
  \qquad
  \Psi=L(t)P.
\end{equation}
完整物理 IPM 严格化为
\begin{equation}
  (U+cz-be_1)\cdot\nabla R=0,
  \qquad
  -\Delta P=R_X,
  \qquad U=\nabla^\perp P,
  \qquad P(X,0)=0,
  \label{eq:beta-zero-bvp}
\end{equation}
其中 $a'(t)=b$。若~\eqref{eq:beta-zero-bvp} 存在有界光滑解且
$R_X\in C_0$、$R_X(z_*)\ne0$，则
\begin{equation}
  \rho_x=\frac{R_X}{c(T-t)},
  \qquad
  \norm{\rho_x}_\infty
  =\frac{\norm{R_X}_\infty}{c(T-t)},
  \label{eq:beta-zero-type-i}
\end{equation}
给出所需的精确 Type-I 爆破。

在 wall transition 上写 $r_0'(0)=\alpha$ 并去掉平移 gauge，方程逐阶强制
\begin{equation}
  R=\alpha X-\frac{\alpha^2}{2c}Y+\Order(|z|^2),
  \qquad
  P=cXY-\frac{\alpha}{2}Y^2+\Order(|z|^3).
  \label{eq:beta-zero-inner-jet}
\end{equation}
所以 affine 核仍是正确内端，但不能延伸到无穷远。另一方面设
\begin{equation}
  R=F_0(\theta)+r^{-1}F_1(\theta)+\cdots,
  \qquad
  P=rG_0(\theta)+G_1(\theta)+\cdots,
\end{equation}
则首阶 Dirichlet--Poisson 方程与 Fredholm 条件为
\begin{equation}
  G_0''+G_0=\sin\theta F_0',
  \qquad
  \int_0^\pi\sin^2\theta F_0'(\theta)\,d\theta=0.
\end{equation}
一个完全显式的非平凡相容外端是
\begin{equation}
  \boxed{
  F_0=A\left(\cos\theta+\frac53\cos3\theta\right),
  \qquad
  G_0=-\frac{4A}{3}\sin^4\theta,}
  \label{eq:beta-zero-tail}
\end{equation}
并有 $R_X=\Order(r^{-1})$、$F_1=G_0F_0'/c$。这证明方向常数远尾与
Poisson、壁面和第一次输运修正之间没有 leading-order 矛盾。尚未完成的是
把~\eqref{eq:beta-zero-inner-jet} 与~\eqref{eq:beta-zero-tail} 连成真正二维、
无限角模的异宿。

而且 leading far end 可以真正集中到任意给定的右侧锐角扇。任取
$0<\theta_-<\theta_+<\pi/2$ 与
$0\not\equiv G_0\in C_c^\infty((\theta_-,\theta_+))$，定义
\begin{equation}
  \boxed{F_0'(\theta)=\frac{G_0''(\theta)+G_0(\theta)}{\sin\theta}.}
  \label{eq:compact-acute-tail}
\end{equation}
则 leading Poisson 方程按定义成立，而 Fredholm 条件由
\begin{equation}
  \int_0^\pi\sin^2\theta F_0'\,d\theta
  =\int_0^\pi\sin\theta(G_0''+G_0)\,d\theta=0
\end{equation}
自动满足。因此
$R_X=-r^{-1}\sin\theta F_0'(\theta)+\Order(r^{-2})$ 只在指定锐角扇内
active。若再要求扇形两侧 $F_0$ 取同一常数，只需用两个 compact bumps 的
线性组合满足额外矩条件 $\int F_0'=0$。具体地，
\begin{equation}
  \int F_0'\,d\theta
  =2\int G_0(\theta)\csc^3\theta\,d\theta,
\end{equation}
故可取两个不交 bumps $\phi_1,\phi_2$，令
$G_0=\phi_1-(\int\phi_1\csc^3/\int\phi_2\csc^3)\phi_2$。非零
$F_0'$ 必须变号；也就是说
right-fan localization 本身就要求 signed angular return。这证明“集中在右侧
锐角扇 + 全方向 $R_X\to0$”在 leading data 层面完全相容。

脚本 \texttt{analysis/ipm\_plot\_compact\_acute\_tail.m} 对一个
$8^\circ<\theta<38^\circ$ 的 two-bump 例作了独立审计：leading Poisson、
Fredholm、same-constant return 残差分别为
$2.22\times10^{-16}$、$1.50\times10^{-18}$、$5.80\times10^{-18}$，扇外
activity 精确为零。图~\ref{fig:compact-acute-tail} 画的是 asymptotic data，
不是完整 nonlinear fixed profile。
\begin{figure}[htbp]
  \centering
  \resultfigure[width=0.96\linewidth]{compact_acute_beta0_tail.png}
  \caption{右侧锐角扇内的 compact-angular leading tail。左列显示
  $F_0,G_0,F_0'$ 与 $rR_X$；右图显示 $R_X=-r^{-1}\sin\theta F_0'$。
  正负 angular return 保证扇形两侧密度常数相同。}
  \label{fig:compact-acute-tail}
\end{figure}

壁面条件事实上给出更强的逐点恒等式
\begin{equation}
  \bigl(cX-b-P_Y(X,0)\bigr)R_X(X,0)=0.
  \label{eq:beta-zero-wall-product}
\end{equation}
若 $P=rG_0(\theta)+o_{C^1}(r)$，则壁面上 $P_Y=\Order(1)$；因此
$R_X(X,0)$ 在两条充分远的 wall tails 上必须\emph{恒等于零}，而不只是趋零。
非平凡候选只能有端点无限阶平坦的 $C^\infty$ compact active segment；实解析
wall trace 被~\eqref{eq:beta-zero-wall-product} 排除。若令
$a(X)=R_X(X,0)$，逐阶展开还给出
\begin{equation}
  p_2=-\frac a2,\qquad r_1=-\frac{a^2}{2c},\qquad
  p_3=\frac{aa'}{6c},\qquad
  r_2=\frac{5a^2a'}{16c^2}.
  \label{eq:beta-zero-wall-hierarchy}
\end{equation}
所以只要 $a$ 不是常数，就必然产生无限 transverse hierarchy；affine 核只是
常斜率时的截断特例。用壁面足点 $q$ 标记特征，$X=X(q,Y)$、$R=F(q)$，输运
可精确积分为
\begin{equation}
  P(q,Y)=c\left(2\int_0^YX(q,s)\,ds-YX(q,Y)\right)-bY,
  \label{eq:beta-zero-hodograph}
\end{equation}
记 $J=X_q>0$、$K=X_Y$、$I=\int_0^YX(q,s)\,ds$、
$\mathcal Q=2I_q-YJ$、$H=X-YK$，则 Poisson 精确等价于单个标量方程
\begin{multline}
 -\partial_q\left[c\frac{1+K^2}{J}\mathcal Q-K(cH-b)\right]\\
 -\partial_Y\left[-cK\mathcal Q+J(cH-b)\right]=F'(q).
 \label{eq:beta-zero-hodograph-pde}
\end{multline}
有效特征场还有精确公式
\begin{equation}
 W=U+cz-be_1=\frac{2cI_q}{J}(K,1).
 \label{eq:beta-zero-hodograph-field}
\end{equation}
故 $J>0$ 时 active 特征向上开放。壁面上 Poisson 强制
$X_Y(q,0)=F'(q)/(2c)$、
$X_{YY}(q,0)=F'(q)F''(q)/(8c^2)$。这严格排除了直 ridge 和共同宽度
$X/s(Y)$；所求特征族必须从 wall 起就弯曲。这里 $F'$ 是 PDE 的自然输出，
不能再作为独立 Neumann 数据。以 $X(q,Y)$ 和 top outer ray $\xi(q)$ 为未知时，
interior PDE、bottom value、两侧 exterior DtN 与 top value/derivative
polyhomogeneous matching 的离散条件数恰等于未知量数，故正确的自由边界
Newton 系统并不过定。

更关键的是，\eqref{eq:beta-zero-tail} 不能在纯 Laurent 类中继续。若写
\begin{equation}
  R\sim\sum_{n\ge0}r^{-n}F_n(\theta),\qquad
  P\sim\sum_{n\ge0}r^{1-n}G_n(\theta),
\end{equation}
则
\begin{equation}
  G_n''+(n-1)^2G_n
  =n\cos\theta F_n+\sin\theta F_n'.
\end{equation}
在 $n=3$，右端对 Dirichlet kernel $\sin2\theta$ 的投影严格为
\begin{equation}
  [S_3]_{\sin2\theta}
  =-\frac{25}{48}\frac{A^4}{c^3}\ne0.
  \label{eq:beta-zero-third-resonance}
\end{equation}
前一级自由 harmonic mode 不能改变该系数。由
$-\Delta(r^{-2}\log r\sin2\theta)=4r^{-4}\sin2\theta$，最小修正被强制为
\begin{equation}
  \boxed{P_{3,\log}
  =\frac{25}{192}\frac{A^4}{c^3}
  r^{-2}\log r\,\sin2\theta.}
  \label{eq:beta-zero-forced-log}
\end{equation}
随后输运还强制 $R$ 出现 $r^{-4}\log r$ 项。因此正确的外端是
polyhomogeneous log cascade；指数小修正位于所有代数阶之外，不能消除
\eqref{eq:beta-zero-third-resonance}。此前提出的对数修正在这里不是装饰，而是
方程必需项。

该共振也不是显式三角例的偶然。对一般 $h=F_0'$、$b=0$，前三阶输运递推为
\begin{gather}
 F_1=\frac{G_0h}{c},\qquad
 F_2=\frac{(G_0F_1)'}{2c},\\
 G_1''=(\sin\theta F_1)',\qquad
 G_2''+G_2=2\cos\theta F_2+\sin\theta F_2',\\
 F_3=\frac{2G_0'F_2+G_0F_2'+G_1'F_1-G_2h}{3c}.
 \label{eq:beta-zero-general-recursion}
\end{gather}
若 $S_n=n\cos\theta F_n+\sin\theta F_n'$，分部积分严格给出
\begin{equation}
 \left\langle S_n,\sin((n-1)\theta)\right\rangle
 =(n-1)\int_0^\pi\sin((n-2)\theta)F_n(\theta)\,d\theta.
 \label{eq:beta-zero-resonance-functional}
\end{equation}
因此 $n=2$ 自动相容，首个可能障碍为 $2\int\sin\theta F_3$；若它非零，
强制项是
\begin{equation}
 \boxed{P_{3,\log}=-\frac1\pi
 \left(\int_0^\pi\sin\theta F_3\,d\theta\right)
 r^{-2}\log r\,\sin2\theta.}
 \label{eq:beta-zero-general-log}
\end{equation}
$G_2$ 的自由 $\lambda\sin\theta$ 模不改变这个泛函，因为其改变量正比于
leading Fredholm 矩。对一个 two-bump compact-right-fan 例，独立 $2401$ 点
sine audit 给出的 log 系数约为 $4.64$，且随加密稳定；这是该 compact acute
子族中共振非零的数值证据，不是全族严格定理。

形式递推还表明，每个 $F_n$ 和 density log 系数的角支撑仍留在原锐扇内；
但 $G_n$ 由全局角 ODE 求出，共振 harmonic 也传播到所有角度。所以目标可以
同时具有“$\rho_x$ 锐扇集中”和“$\rho_x\to0$”，却一般不能再要求 $P,U$
也保持角紧支撑。

第一轮有限维搜索也给出了有用的否证。脚本
\texttt{analysis/ipm\_search\_beta0\_profile.m} 精确消去离散 Poisson 方程，
只在输运入流边给密度数据，并弱匹配外圈 $R_X\to0$。在 $129\times49$ 网格、
$32\times14$ 模态下，Poisson 残差为 $1.85\times10^{-11}$，但 bulk/wall
输运 RMS 仍为 $0.2104/0.1044$，外圈 $\max|R_X|=0.2424$，且 wall return
出现高频振荡。因此该结果明确是 noncandidate；下一轮应直接离散
\eqref{eq:beta-zero-hodograph} 和强制 log outer data，而不是继续给直模加阶。

\subsubsection{标准 Green 自屏蔽与动态 Casimir cap}

若不允许另加自由的 $BXY$ harmonic 场，局部双曲应变必须由 density 自己
产生。半平面 Dirichlet Green 核给出精确本征值条件
\begin{equation}
 \boxed{c=P_{XY}(0,0)=\frac2\pi\operatorname{p.v.}
 \int_{\mathbb H}\frac{XY}{(X^2+Y^2)^2}R_X\,dX\,dY.}
 \label{eq:self-screened-moment}
\end{equation}
可分部积分时也可写成
\begin{equation}
 P_{XY}(0,0)=\frac2\pi\operatorname{p.v.}\int_{\mathbb H}
 R\frac{Y(3X^2-Y^2)}{(X^2+Y^2)^3}\,dX\,dY.
\end{equation}
若 bounded return 全在重整化距离 $r\ge D$，则
\begin{equation}
 |P_{XY}^{\rm far}(0,0)|\le\frac{4\norm{R}_\infty}{\pi D},
 \qquad
 R\in L^1\Longrightarrow
 |P_{XY}^{\rm far}(0,0)|\le\frac{2\norm{R}_1}{\pi D^3}.
 \label{eq:self-screened-far-bound}
\end{equation}
所以外逃 reservoir 只能保存 Casimir 和截断远场，不能单独产生 order-one
inner strain；signed head 必须留在重整化 $O(1)$ 区域。

这个 head 在椭圆层面有显式紧支撑构造。取
\begin{equation}
 Q=R_X=X\phi(r),\qquad \phi\in C_c^\infty((r_1,r_2)),
 \qquad R(X,Y)=\int_{-\infty}^XQ(S,Y)\,dS.
 \label{eq:self-screened-annulus}
\end{equation}
则每条水平线上 $\int Q\,dX=0$，$R,Q$ 光滑紧支撑，且
\begin{equation}
 \boxed{P_{XY}(0,0)=\frac4{3\pi}\int_0^\infty\phi(r)\,dr.}
 \label{eq:self-screened-annulus-strain}
\end{equation}
令 $\int\phi=3\pi c/4$ 就得到 $P=cXY+\Order(r^4)$。右侧正 $Q$ 与左侧
负 return 的符号正是 source-neutral 所需；再加角权可把正梯度集中到右锐扇。
因此 Poisson、多极矩和正负号都不是障碍，但
\eqref{eq:self-screened-annulus} 只是一时刻的椭圆 gluing，不是动态解。

事实上 $W=U+cz-be_1$ 满足 $\nabla\cdot W=2c$。若 stationary profile 的
每个非背景 level set 都有有限面积，对支撑于非背景值域的非负 $h(R)$ 有
\begin{equation}
 0=\int_{\mathbb H}\nabla\cdot(Wh(R))
 =2c\int_{\mathbb H}h(R),
 \label{eq:beta-zero-compact-nogo}
\end{equation}
矛盾。所以非平凡 compact/finite-level-area 的 stationary $\beta=0$ profile
不存在；必须保留方向常数无限面积尾，或转向非定常 return。

后一条路线的面积匹配是严格相容的。令 $s=-\log(T-t)$，则
\begin{equation}
 cR_s+(U+cz-be_1)\cdot\nabla R=0,
 \qquad -\Delta P=R_X.
 \label{eq:beta-zero-nonstationary}
\end{equation}
若物理初值的 $\lambda$-超水平集面积为 $A_\lambda$，而局部角尾在
$r_z\sim e^sa(\theta)$ 处接 cap，则不可压缩性对每个 level 强制
\begin{equation}
 \boxed{A_\lambda=\frac{c^2}{2}
 \int_{\{F_0(\theta)>\lambda\}}a(\theta)^2\,d\theta.}
 \label{eq:beta-zero-area-law}
\end{equation}
cap 的物理半径为 $O(1)$，而它对 inner $P_{XY}$ 的影响只有
$O(e^{-s})=O(T-t)$。于是“局部 signed head 自生 $1/(T-t)$ strain + 固定物理
尺度 cap 保存全部 Casimir”在尺度和椭圆矩上完全相容；真正未解的是满足
\eqref{eq:beta-zero-nonstationary} 的动态连接层。

这个更新也改变了 Bessel 的定位：有限 $L^p$ Casimir 和势能排除不再能用于
方向常数尾，但 local 多极矩、wall $D$-层和有限 transverse-mode 障碍仍然
有效。因此 Bessel 应作为中间匹配环带的 Green corrector，而最终外端由
angular-constant/signed-return 层承担。完整审计另见
\texttt{analysis/RHOX\_DECAY\_BLOWUP\_STATUS\_ZH.md}。

\subsection{开放 harmonic contract 下的解析 Type-I 精确轨道}

若只要求~\eqref{eq:rhox-c0}，而不要求速度采用衰减 Biot--Savart 归一化，则
可以构造一条密度初态光滑、但远场合同预先含奇异时刻 $T$ 的精确轨道。令
$\tau=T-t$、
$\sigma=\tau^2/2$、$X=x/\tau$、$Y=\tau y$，并取
\begin{equation}
  \Psi_H=\frac{xy}{\tau},\qquad
  \Psi_s=\tau\Phi(\sigma,X,Y),\qquad
  \rho=\rho_c+F(\sigma,X,Y).
\end{equation}
完整无源 IPM 精确等价于
\begin{equation}
  F_\sigma=\{\Phi,F\},\qquad
  -(\Phi_{XX}+4\sigma^2\Phi_{YY})=F_X,
  \qquad \rho_x=\tau^{-1}F_X.
  \label{eq:affine-contract-profile}
\end{equation}
对 $Y$ 作 Fourier 变换，椭圆逆的精确核为
\begin{equation}
  \widehat\Phi(X,\eta)
  =-\frac12\int_{\mathbb R}\operatorname{sgn}(X-Z)
  e^{-2\sigma|\eta||X-Z|}\widehat F(Z,\eta)\,dZ.
  \label{eq:affine-contract-kernel}
\end{equation}
它把 $U$ 的 $L_X^1$ 解析范数一致映到无权 $L_X^\infty$ 速度范数。衰减权重
应放在被输运的 $U=\nabla F$ 上，不能放到具有 Poisson 远尾的
$\mathcal K_\sigma U$ 上。具体取
\begin{equation}
  \norm{U}_r=\sum_\alpha\frac{r^{|\alpha|}}{\alpha!}
  \left(\norm{\partial^\alpha U}_\infty
  +\norm{\partial^\alpha U}_{L_X^1L_Y^\infty}\right),
\end{equation}
则对 $0<r'<r$，一致于 $\sigma\ge0$，
\begin{equation}
  \norm{\mathcal N_\sigma(U)-\mathcal N_\sigma(V)}_{r'}
  \le\frac{C}{r-r'}(\norm U_r+\norm V_r)\norm{U-V}_r.
\end{equation}
标准 Ovsyannikov/Cauchy--Kowalevski 迭代由此从 $\sigma=0$ 给出局部唯一解析
衰减解。Gaussian profile 速度及其所有导数有界，特征流与恒等映射相差有界且
各阶导数有界，故 Schwartz 密度扰动由与逆流复合而保持 Schwartz。

全平面取
\begin{equation}
  F_*(X,Y)=\varepsilon e^{-X^2-Y^2}.
\end{equation}
取 $\rho_c=0$，将该终端剖面向 $0\le\sigma\le\sigma_0$ 延拓并按
\eqref{eq:affine-contract-profile} 拉回，得到解析 Schwartz 密度扰动初态的精确
轨道，且
\begin{equation}
  \rho_x(t,\cdot)\in C_0(\mathbb R^2),\qquad
  \tau\norm{\rho_x(t)}_\infty
  \longrightarrow |\varepsilon|\sqrt{\frac2e}>0.
  \label{eq:affine-contract-blowup}
\end{equation}
峰位满足 $(x,y)=(\pm\tau/\sqrt2+o(\tau),O(\tau))$，所以不向无穷远逃逸。
但是它也不是孤立点集中：对每个固定有限 $y_0$，同一 $x=O(\tau)$ 层内仍有
$|\rho_x|\asymp\tau^{-1}$；更精确地，
\begin{equation}
  \tau\rho_x(t,\pm\tau/\sqrt2,y_0)
  \longrightarrow\mp\varepsilon\sqrt{2/e}.
\end{equation}
物理解在 $x$ 方向收缩、在 $y$ 方向展宽，极限
奇异集是竖直 sheet，而不是右侧锐角点。它是 literal $\rho_x\in C_0$ 条件下
的 Type-I 精确轨道，但不能代替最终锐角目标。半平面 odd 数据
$F_*=\varepsilon e^{-X^2}Ye^{-Y^2}$ 也给
$\tau\norm{\rho_x}_\infty\to|\varepsilon|/e$，但峰位
$y\sim(\sqrt2\tau)^{-1}$ 逃向无穷远，不能算固定壁点奇点。

这个正例的合同边界必须写清：$\Psi_H=xy/\tau$ 二次增长、其速度线性增长、
无限能且
$\norm{\nabla u_H}\sim\tau^{-1}$。它逐点没有体源，却要求预先给定
$M(t)=\operatorname{diag}(-1/\tau,1/\tau)$ 的时间奇异远场压力/速度合同；这个
未来合同不由 $t=0$ 的光滑密度 Cauchy 数据决定。因此它不是衰减 Dirichlet
Biot--Savart 唯一选出的自治 Cauchy 正例；真正锐角 gluing 仍需用 density
return 自洽地产生并屏蔽这个局部外应变。

这个替换必然是非微扰的。若去掉 $\Psi_H$，仍用同一面积保持坐标和
$\Psi_s=\tau\Phi$，则
\begin{equation}
  2\sigma F_\sigma
  =XF_X-YF_Y+2\sigma\{\Phi,F\}.
  \label{eq:self-screened-affine}
\end{equation}
若 $F\to F_*$ 于局部 $C^1$ 且 self bracket 有界，终端方程为
$XF_{*X}-YF_{*Y}=0$，故 $F_*=H(XY)$。而
$F_{*X}=YH'(XY)\in C_0$ 强制 $H'\equiv0$：固定 $q=XY$，令
$Y\to\infty$、$X=q/Y$ 即可。因此任何非平凡 self-screened return 都必须使
$\{\Phi,F\}=\Order(\sigma^{-1})$，由 wall fast layer、锐角内核或进入 leading
order 的 return 激活退化椭圆坏频带；不能把上面的 harmonic trajectory 当作
标准 Cauchy 解的小修正。

\subsubsection{壁面 Hilbert 闭合与动态 edge 的必要性}

同一面积保持坐标还能严格审计 density 能否自生取代 $\Psi_H$。令
$s=-\log\tau$，去掉外加 harmonic strain 后，完整方程是
\begin{equation}
 F_s+XF_X-YF_Y+\tau^2\{\Phi,F\}=0,
 \qquad
 -(\Phi_{XX}+\tau^4\Phi_{YY})=F_X.
 \label{eq:wall-hilbert-system}
\end{equation}
半平面 Dirichlet--Neumann 数据满足
\begin{equation}
 \widehat{\Phi_Y}(k,0)=\frac{ik}{\tau^4}
 \int_0^\infty e^{-|k|Y/\tau^2}\widehat F(k,Y)\,dY.
 \label{eq:wall-hilbert-neumann}
\end{equation}
若 wall fast layer 中 $F\to f(X)$ 足够光滑，则
$\tau^2\Phi_Y(X,0)\to-\mathcal Hf(X)$，所以 stationary wall trace 必须满足
\begin{equation}
 \boxed{(X+\mathcal Hf)f'=0.}
 \label{eq:wall-hilbert-closure}
\end{equation}
在单一 active interval $(-a,a)$ 且两端接同一常数时，有限 Hilbert 反演给出
\begin{equation}
 f(X)=C-\sqrt{a^2-X^2},\qquad |X|<a,
 \label{eq:wall-hilbert-root}
\end{equation}
因为 $\mathcal H[\sqrt{a^2-\cdot^2}_+]=X$。然而
$f'=X/\sqrt{a^2-X^2}$ 在端点强制平方根奇性。因此单区间、完全 stationary、
自屏蔽的 wall closure 不可能同时有 $C^\infty$ compact-active edge；从光滑
初值出发必须让 edge/return 随 $s$ 运动，形成真正的第二尺度。

局部 strain 符号没有问题：
\begin{equation}
 \lim_{\tau\to0}\tau\Psi_{xy}(0,0)=-\Lambda f(0).
 \label{eq:wall-strain-selection}
\end{equation}
例如 $f=-\varepsilon e^{-X^2}$ 给
$-\Lambda f(0)=2\varepsilon/\sqrt\pi>0$，只是在整个 active interval 上并不
满足 $\mathcal Hf=-X$。这说明最终对象应是 local signed head、动态平方根
edge regularization 与 $r_z\sim e^s$ Casimir cap 的三层连接，而不是另一个
全局光滑静态 wall fixed point。

\section{进一步可得的结果与下一步}

\subsection{极限热方程的条件稳定性}

对 $K_0$ 极限写
\begin{equation}
  \mathcal R(X,\eta)=X^{-1/2}G(s,\zeta),
  \qquad s=\log X,
  \qquad \zeta=\eta/\sqrt X.
\end{equation}
热方程变成 Ornstein--Uhlenbeck 型方程
\begin{equation}
  G_s=G_{\zeta\zeta}+\frac12\zeta G_\zeta+\frac12G.
  \label{eq:ou}
\end{equation}
Gaussian 是固定点；在合适 Gaussian 加权空间中，质量模态为中性，平移、
方差及更高 Hermite 模态依次衰减。由此可以得到一个\emph{极限线性问题内}
的条件稳定结论，并明确数值中要调制掉的中性参数：幅度、法向平移 $h$ 和
纵向尺度。它尚不是完整 IPM 的稳定性定理，因为~\eqref{eq:degenerate-limit}
中的一般退化 IPM 输运耦合仍是 $\Order(1)$；只有在特殊闭包 $P=-R$ 上才
退化成热方程。此外还需统一控制 $\Order(\delta)$ 的纵向椭圆修正、壁面
条件和内外区匹配。

\subsection{建议的四阶段研究计划}

\paragraph{阶段 A：直接解自洽 hodograph 自由边界。}
以 $X(q,Y)$ 和 outer ray $\xi(q)$ 为 unknown，直接离散
\eqref{eq:beta-zero-hodograph-pde}。wall basis 写成端点无限阶平坦的 compact
$k(q)=X_Y(q,0)$，自动编码 $X_{YY}=kk'/2$；两侧只施加 exterior harmonic
DtN，top 匹配含 \eqref{eq:beta-zero-general-log} 的 polyhomogeneous radiation。
每次 Newton 同时审计 \eqref{eq:self-screened-moment}，不能把 $c$ 当外加参数。

\paragraph{阶段 B：构造非定常 signed-head/edge/cap 连接。}
以 $Q=R_X\in C_0$ 为主变量，从
\eqref{eq:self-screened-annulus} 的 source-neutral signed head 出发，演化
\eqref{eq:beta-zero-nonstationary}；用 \eqref{eq:beta-zero-area-law} 锁定每个
level 的外 cap。重点解析 \eqref{eq:wall-hilbert-root} 的平方根端点如何在一个
随 $s$ 运动的 fast edge 中被光滑化，并验证 head 贡献保持 order one、cap
贡献按 $e^{-s}$ 衰减。

\paragraph{阶段 C：从无预置奇核的光滑初值反向验证。}
在完整物理 IPM 中使用光滑紧支撑或快速衰减数据，独立测量 head 尺度、edge
厚度和 cap 半径；检验 $\norm{\rho_x}_\infty\sim(T-t)^{-1}$、固定右锐扇内的
归一化 $\rho_x$、全方向远端衰减、面积律和 self-strain moment 是否同时平台化。
任何结论都必须通过网格、盒子、时间步和初值族四重比较。

\paragraph{阶段 D：只把 Bessel 当作中间 Green corrector。}
对解析 $K_0/K_1$ 家族保留 $\varepsilon$ 扫描和零线回归，用它们检验局部
Poisson 响应；但不再要求物理解收敛到一个 Bessel fixed point。只有当上述
signed-head 到 angular/log tail 的连接产生自然中间尺度时，才把 Bessel
匹配残差投影到平移、幅度和缩放零模并推导真实调制方程。

\subsection{可以被证伪的判据}

一个可信的“光滑初值选择两尺度稳态”结论至少应同时满足：
\begin{enumerate}
  \item 两个物理长度都按同一候选奇异时刻收缩，并且比例律在加密网格后稳定；
  \item 归一化二维场、$\rho_x$ 零线和完整 PDE 残差同时收敛；
  \item 去掉初值中预置的 Bessel 长尾或改变光滑核心后仍选择同一轮廓；
  \item 扩大物理盒子后，中心动力学不由开放边界或 Green 截断决定；
  \item 最大值原理、质量/通量预算与物理坐标回代均闭合；
  \item 与 Lipschitz 锐角爆破、无限能特殊 ansatz、带外源爆破明确区分。
\end{enumerate}

\section{结论}

穿壁 $K_0$ 与齐次壁面 $K_1\sin\theta$ 仍给出精确 Bessel/heat 两尺度形状，
其 $\rho_x$ 零线分别为 $\sqrt{2X}$ 与 $\sqrt{6X}$；但它们描述的是严格
抛物远场极限，不证明物理解从光滑初值自动收缩到尖点。Bessel 在最终构造中
应降为可精确审计的中间 Green corrector。

在全方向 $\rho_x\in C_0$ 的硬条件下，现在已经得到三个更直接的结果。
第一，任意指定右侧锐角扇都存在 compact-angular、same-constant-return 的
leading tail，$R_X=O(r^{-1})$ 且扇外为零；高阶 density activity 形式上仍留在
原扇。第二，纯幂展开在首个非平凡 Fredholm 投影处强制
$r^{-2}\log r\sin2\theta$ 流函数项，所以指数修正不能代替对数级联。第三，
标准 Green 核下的 self-strain moment 可由光滑紧支撑 signed head 精确实现，
而 $r_z\sim e^s$ 的外 cap 可同时保存全部 level-area Casimir 且只给 inner
strain 一个 $O(e^{-s})$ 修正。因此“右锐扇 + 远端衰减 + 自生双曲应变”在
椭圆、矩和尺度层面没有矛盾。

剩余障碍已经定位在输运 gluing：finite-level-area 的 compact stationary
$\beta=0$ profile 被散度恒等式排除，单 active interval 的 stationary wall
closure 又被 \eqref{eq:wall-hilbert-root} 的平方根端点排除。因此从全局光滑
半平面初值出发的候选必须是非定常三层对象：重整化 $O(1)$ 的 local signed
head 产生 $1/(T-t)$ strain，动态 fast edge 平滑平方根端点，固定物理尺度的
angular cap 承担 return/Casimir；其外展开含强制 log cascade。下一次计算应
直接求这个 hodograph/return 系统，而不是继续增加旧有限模 ansatz。

严格边界仍需明确：Dembski 的紧支撑锐角域解给出了远端衰减的有限时梯度
爆破，但角点初值一般只到 Lipschitz；\eqref{eq:affine-contract-blowup} 给出解析
Schwartz 密度初态的 exact Type-I 轨道，却预设了未来奇异 harmonic far-field
strain。标准衰减 Biot--Savart、全局光滑半平面初值、固定有限右锐角点的自治
爆破尚未证明；目前最短且可判定的缺口就是上述动态 signed-head/edge/cap
连接。

\begin{thebibliography}{99}

\bibitem{castro2007}
D. C\'ordoba, F. Gancedo, and R. Orive,
``Analytical behavior of two-dimensional incompressible flow in porous media,''
\emph{Journal of Mathematical Physics} 48 (2007), 065206.
\href{https://doi.org/10.1063/1.2404593}{doi:10.1063/1.2404593}.

\bibitem{elgindi2017}
T. M. Elgindi,
``On the asymptotic stability of stationary solutions of the inviscid
incompressible porous medium equation,''
\emph{Archive for Rational Mechanics and Analysis} 225 (2017), 573--599.
\href{https://doi.org/10.1007/s00205-017-1090-7}{doi:10.1007/s00205-017-1090-7}.

\bibitem{kiselevyao2023}
A. Kiselev and Y. Yao,
``Small scale formations in the incompressible porous media equation,''
\emph{Archive for Rational Mechanics and Analysis} 247 (2023), Article 1.
\href{https://doi.org/10.1007/s00205-022-01830-z}{doi:10.1007/s00205-022-01830-z}.

\bibitem{bianchinicordoba2025}
R. Bianchini, D. C\'ordoba, and L. Mart\'inez-Zoroa,
``Non existence and strong ill-posedness in $H^2$ for the stable IPM equation,''
\emph{Journal of Functional Analysis} 289 (2025), 111097.
\href{https://doi.org/10.1016/j.jfa.2025.111097}{doi:10.1016/j.jfa.2025.111097}.

\bibitem{bianchinijoparkwang2025}
R. Bianchini, M. J. Jo, J. Park, and S. Wang,
``Sharp asymptotic stability of the incompressible porous media equation,''
arXiv:2505.05165 (2025).
\url{https://arxiv.org/abs/2505.05165}.

\bibitem{kiselevsarsam2025}
A. Kiselev and N. A. Sarsam,
``Finite time blow-up in a 1D model of the incompressible porous media equation,''
\emph{Nonlinearity} 38 (2025), 055018.
\href{https://doi.org/10.1088/1361-6544/adcdb8}{doi:10.1088/1361-6544/adcdb8}.

\bibitem{cordobamartinez2024}
D. C\'ordoba and L. Mart\'inez-Zoroa,
``Finite time singularities of smooth solutions for the 2D incompressible
porous media (IPM) equation with a smooth source,'' arXiv:2410.22920 (2024).
\url{https://arxiv.org/abs/2410.22920}.

\bibitem{collotprangetan2026}
C. Collot, C. Prange, and J. Tan,
``Stable self-similar singularity formation for infinite energy solutions of
the incompressible porous medium equations,''
\emph{Communications in Mathematical Physics} 407 (2026), Article 150.
\href{https://doi.org/10.1007/s00220-026-05688-0}{doi:10.1007/s00220-026-05688-0}.

\bibitem{dembski2025}
K. H. Dembski,
``Singularity formation in the incompressible porous medium equation without
boundary mass,'' arXiv:2511.01827 (2025; manuscript updated 2026).
\url{https://arxiv.org/abs/2511.01827}.

\bibitem{constantinlavicol2019}
P. Constantin, J. La, and V. Vicol,
``Remarks on a paper by Gavrilov: Grad--Shafranov equations, steady solutions
of the three dimensional incompressible Euler equations with compactly
supported velocities, and applications,''
\emph{Geometric and Functional Analysis} 29 (2019).
\href{https://doi.org/10.1007/s00039-019-00516-1}{doi:10.1007/s00039-019-00516-1}.

\bibitem{wang2025}
Y. Wang et al.,
``Discovery of unstable singularities,'' arXiv:2509.14185 (2025).
\url{https://arxiv.org/abs/2509.14185}.

\bibitem{wangleger2025}
Y. Wang, T. L\'eger, C.-Y. Lai, and T. Buckmaster,
``Resolving sharp gradients of unstable singularities to machine precision
via neural networks,'' arXiv:2511.22819 (2025).
\url{https://arxiv.org/abs/2511.22819}.

\bibitem{dlmf}
NIST Digital Library of Mathematical Functions,
Chapter 10, ``Bessel Functions,'' especially the modified-Bessel asymptotics.
\url{https://dlmf.nist.gov/10}.

\end{thebibliography}

\end{document}
